\begin{frame}{``모델을 안다''의 두 갈래}
  Dyna-Q의 모델 $\hat{P},\hat{R}$은 \textbf{경험에서 쌓인} 것이었다
  (각 $(s,a)$마다 ``마지막으로 관찰된 $(r,s')$''을 기억하는 표).
  \vskip0.6em
  \begin{columns}[T]
    \begin{column}{0.48\textwidth}
      \kb{known model(알려진 모델)}
      \begin{itemize}
        \item 바둑·체스: 모델은 학습된 게 아니라
          \textbf{규칙 자체} --- 게임 규칙을 코드로 옮긴 것
        \item 데이터로부터 추정한 근사가 아니므로
          \textbf{오차가 전혀 없는 완벽한 모델}
      \end{itemize}
    \end{column}
    \begin{column}{0.48\textwidth}
      \kb{learned model(학습된 모델)}
      \begin{itemize}
        \item Dyna-Q처럼 \textbf{데이터로부터 학습}한 모델
        \item 로봇 같은 실제 환경은 보통 이쪽 --- 모델이
          데이터로부터 배워야 하거나(또는 시뮬레이터로
          만들거나) 아예 만들지 못한다
      \end{itemize}
    \end{column}
  \end{columns}
  \vskip0.7em
  \kq{모델이 \textbf{어디에서 오느냐}(known/learned)와
  모델을 \textbf{무엇에 쓰느냐}(학습/계획)는 두 개의 독립된 축이다.}
\end{frame}
